%%%%%%%%%%%%%%%%%%%%%%%%%%%%%%%%%%%%%%%%
%12pt: grandezza carattere
%a4paper: formato a4
%openright: apre i capitoli a destra
%twoside: serve per fare un documento fronteretro
%report: stile tesi (oppure book)
\documentclass[12pt,a4paper,openright,twoside]{report}

\usepackage{tabularx}
\usepackage[english,provide=*]{babel}
\usepackage[utf8]{inputenc}
\usepackage{fancyhdr}
\usepackage{indentfirst}
\usepackage{graphicx}
\usepackage{newlfont}
%librerie matematiche
\usepackage{amssymb}
\usepackage{amsmath}
\usepackage{latexsym}
\usepackage{amsthm}
\usepackage{listings}
\usepackage{chngcntr}
\usepackage[outputdir=build,chapter]{minted}
\usepackage{algpseudocodex}

\oddsidemargin=30pt \evensidemargin=20pt%impostano i margini
\hyphenation{sil-la-ba-zio-ne pa-ren-te-si}

\usepackage[square, numbers, comma, sort&compress]{natbib} %Bibliografia
\usepackage[hidelinks]{hyperref}
\usepackage{float}
\usepackage{caption}
\usepackage[pass]{geometry}

%comandi per l'impostazione della pagina, vedi il manuale della libreria fancyhdr per ulteriori delucidazioni
\captionsetup{font=small}
\setlength{\parindent}{0pt}
\pagestyle{fancy}\addtolength{\headwidth}{20pt}
\renewcommand{\chaptermark}[1]{\markboth{\thechapter.\ #1}{}}
\renewcommand{\sectionmark}[1]{\markright{\thesection \ #1}{}}
\rhead[\fancyplain{}{\bfseries\leftmark}]{\fancyplain{}{\bfseries\thepage}}
\cfoot{}

\linespread{1.3} %comando per impostare l'interlinea

\renewcommand{\listingscaption}{Codice}
\renewcommand{\listoflistingscaption}{List of Algorithms}

\begin{document}
%%%%%%%%%%%%%%%%%%%%%%%%%%%%%%%%%%%%%%%%
% FRONTESPIZIO
\newgeometry{hmarginratio=1:1}
\begin{titlepage}
\begin{center}
\includegraphics[width=2.56in]{figures/logo/logo_unibo.png}\\

\vspace{5mm}
{\small{\bf SCUOLA DI INGEGNERIA E ARCHITETTURA\\
\vspace{2mm}
Dipartimento di Informatica -- Scienza e Ingegneria\\
\vspace{2mm}
Corso di Laurea in Ingegneria Informatica }}
%se Laurea Magistrale scrivere "Corso di Laurea Magistrale in Ingegneria Informatica"
\end{center}
\vspace{11.9mm}
\begin{center}
{\LARGE{\bf \uppercase{Integrating Federated Learning with the Blockchain Infrastructure}}}
\end{center}
\vspace{11.9mm}
\par
\noindent
\begin{minipage}[t]{0.47\textwidth}
{\normalsize{\bf Relatore:\\
Prof. Paolo Bellavista

\vspace{5mm}
}}
\end{minipage}
\hfill
\begin{minipage}[t]{0.47\textwidth}\raggedleft
{\normalsize{\bf Presentata da:\\
Bombarda Alessandro}}
\end{minipage}
\vspace{10mm} % da aumentare a 20 o 30 se non si inseriscono i correlatori
\begin{center}
{\normalsize{\bf Sessione 3 \\ %inserire il numero della sessione in cui ci si laurea
Anno Accademico 2023/2024}} %inserire l'anno accademico a cui si è iscritti
\end{center}
\end{titlepage}

\restoregeometry

%%%%%%%%%%%%%%%%%%%%%%%%%%%%%%%%%%%%%%%%
% ABSTRACT
\clearpage{\pagestyle{empty}\cleardoublepage}
\pagenumbering{roman}
\renewcommand{\abstractname}{Introduction}
\phantomsection
\addcontentsline{toc}{chapter}{Introduction}
\begin{abstract}
  In recent years, advancements in distributed computing and decentralized technologies have reshaped the
  landscape of machine learning and data processing. This document explores the intersection of two
  innovative paradigms: Federated Learning and Blockchain. Each of these fields has gathered
  significant attention for its unique potential — federated learning for enabling collaborative model training
  without centralized data aggregation, and blockchain for providing decentralized governance and immutable ledgers
  with robust security guarantees.
  
  Federated learning is a decentralized approach to training machine learning models collaboratively across
  multiple devices or nodes while keeping data local to each participant. Originally developed to address
  privacy concerns, Federated Learning has found applications in domains such as healthcare, mobile devices, and IoT, where
  sensitive data remains confined to the edge devices that generate it. However, Federated Learning is not without its
  challenges.
  
  Blockchain, on the other hand, is a decentralized ledger technology that ensures secure and immutable
  decentralized data storage and trustless interactions among participants.
  First popularized by cryptocurrencies such as Bitcoin, blockchain has since evolved to support a wide range
  of use cases, including supply chain management, smart
  contracts, and secure data sharing.
  
  The possibility of integrating blockchain with federated learning has sparked interest in the research community as a
  means to address some of the inherent limitations of the paradigm.
  Federated learning's reliance on a central server for model aggregation introduces risks of a single point of failure
  and trust issues among participants.
  Blockchain's decentralized nature and consensus mechanisms can solve this problem among others, such as
  ensuring model integrity and fairness in participation. However this comes with its own set of challenges,
  mostly communication, storage and (depending on the underlying mechanism of the blockchain, also computation) overhead.
  
  Integrating blockchain into federated learning systems presents unique challenges. These include the
  computational and communication overhead introduced by blockchain protocols, ensuring scalability in systems
  with millions of devices, and aligning the asynchronous nature of federated learning with blockchain's
  consensus requirements. Despite these hurdles, the potential benefits—such as enhanced security,
  transparency, and trust—make this integration a promising avenue of research.
  
  This thesis aims to explore the foundational concepts of both federated learning and blockchain, 
  analyze their potential integration, and address the challenges that arise from combining these 
  technologies. Through theoretical insights and practical implementation, the study will evaluate the 
  feasibility and benefits of using blockchain to enhance federated learning systems.
\end{abstract}

%%%%%%%%%%%%%%%%%%%%%%%%%%%%%%%%%%%%%%%%
% RINGRAZIAMENTI
% \clearpage{\pagestyle{empty}\cleardoublepage}
% \chapter*{Ringraziamenti}
% \phantomsection
% \addcontentsline{toc}{chapter}{Ringraziamenti}
% TODO

\clearpage{\pagestyle{empty}\cleardoublepage}

%%%%%%%%%%%%%%%%%%%%%%%%%%%%%%%%%%%%%%%%
% INDICE
\tableofcontents %crea l'indice
%imposta l'intestazione di pagina
\rhead[\fancyplain{}{\bfseries\leftmark}]{\fancyplain{}{\bfseries\thepage}}
\lhead[\fancyplain{}{\bfseries\thepage}]{\fancyplain{}{\bfseries
Index}}

%%%%%%%%%%%%%%%%%%%%%%%%%%%%%%%%%%%%%%%%
% ELENCO DELLE FIGURE
\clearpage{\pagestyle{empty}\cleardoublepage}
\renewcommand{\listfigurename}{List of Figures}
\phantomsection
\addcontentsline{toc}{chapter}{List of Figures}
\listoffigures %crea l'elenco delle figure

%%%%%%%%%%%%%%%%%%%%%%%%%%%%%%%%%%%%%%%%
% ELENCO DELLE TABELLE
\clearpage{\pagestyle{empty}\cleardoublepage}
\renewcommand{\listtablename}{List of Tables}
\phantomsection
\addcontentsline{toc}{chapter}{List of Tables}
\listoftables %crea l'elenco delle tabelle

%%%%%%%%%%%%%%%%%%%%%%%%%%%%%%%%%%%%%%%%
% ELENCO DEI CODICI
\clearpage{\pagestyle{empty}\cleardoublepage}
\phantomsection
\addcontentsline{toc}{chapter}{List of Algorithms}
\listoflistings %crea l'elenco dei codici

%%%%%%%%%%%%%%%%%%%%%%%%%%%%%%%%%%%%%%%%
% INTRODUZIONE
\clearpage{\pagestyle{empty}\cleardoublepage}
\markboth{}{} % Clears the header markers
\pagenumbering{arabic} %mette i numeri arabi
\lhead[\fancyplain{}{\bfseries\thepage}]{\fancyplain{}{\bfseries\rightmark}}



%%%%%%%%%%%%%%%%%%%%%%%%%%%%%%%%%%%%%%%%
% SCENARI APPLICATIVI E STATO DELL'ARTE

\clearpage{\pagestyle{empty}}
\section{Preliminaries}\label{chap:introduction}

\subsection{Context of the Thesis}\label{chap:contextandobjectives}
Advancements in technology have driven a data-centric shift across industries such as healthcare, finance,
smart cities, and the Internet of Things (IoT). This growing reliance on data has introduced significant
challenges in privacy, security, and scalability, necessitating innovative solutions. This thesis explores
the intersection of two emerging paradigms—Federated Learning (FL) and Blockchain technology—as a potential
means to address these challenges.

\textbf{Federated Learning} has emerged as a transformative approach to machine learning by enabling
collaborative model training without the need to share raw data. This is achieved by keeping relevant
computation local to devices and only sharing updates with the interested party.
This preserves data privacy while still achieving high-performance learning. Despite its benefits, FL faces
critical challenges, including reliance on a central server, susceptibility to malicious participants, and a
lack of transparency in the training process.

\textbf{Blockchain technology}, initially developed as the backbone for cryptocurrencies like Bitcoin,
provides a decentralized, secure, and transparent framework for managing data and operations among untrusted
parties. With its core features—such as consensus mechanisms, immutability, and smart contracts—Blockchain
has demonstrated value in sectors ranging from finance to supply chain management. However, its scalability
limitations and resource-intensive operations can hinder its applicability in dynamic environments.

Integrating FL and Blockchain presents a promising avenue for overcoming the limitations of both
technologies. Blockchain's decentralization can eliminate the dependency on a central server in FL, while its
transparency and consensus protocols enhance trust and accountability. Conversely, FL's distributed learning
capabilities reduce the need for raw data storage on the Blockchain, mitigating its scalability concerns.

This thesis is motivated by real-world challenges such as:
\begin{itemize}
  \item Ensuring privacy and security in data-sharing across regulated sectors like healthcare and finance.
  \item Establishing trust in collaborative environments where participants may not fully trust each other,
    such as in decentralized IoT systems.
\end{itemize}

While Federated Learning and Blockchain are an ongoing field of research and a fair amount of articles have
been published about these topics,
their integration has been underexplored. Combining the two technologies introduces novel challenges
including communication privacy and efficiency, synchronism of updates, security and scalability.

Through the design and evaluation of a Blockchain-enhanced Federated Learning framework, this thesis aims to
provide a foundational understanding of the two technologies and explore the challanges and opportunities in
context of privacy, trust, and scalability. The following section outlines the specific objectives of this
thesis, which are derived from this contextual foundation.

\subsection{Objectives of the Thesis}

This document aims to explore the feasibility and practicality of implementing a federated learning system on
a blockchain network. The contributions of this work are threefold:

\paragraph{Explore the technologies}
The thesis seeks to provide to the reader a foundational understanding of both Federated Learning and Blockchain
technologies. We will try to be as thorough as possible, within the limits of the purposes of this document,
highlighting the individual strengths, weaknesses, and potential areas of
application of both technologies.

\paragraph{Exploring the integration}
After providing an understanding of the technologies and their strengths and weaknesses we want to discuss
the advantages and challenges that arise when joining them.

We will discuss them in depth in the coming chapters but as an anticipation, some of the advantages include:
\begin{itemize}
  \item Addressing the central server reliance in FL through Blockchain's decentralized
    infrastructure.
  \item Enhancing trust, transparency, and security in FL systems using Blockchain's
    consensus mechanisms and immutability.
  \item Provide incentives for devices to partake in the federated learning process.
\end{itemize}

And some of the challenges include:
\begin{itemize}
  \item Computational and communication overhead due to Blockchain protocols.
  \item Ensuring security and consistence of the gloabl model.
\end{itemize}

When discussing the challenges and advantages we will also take a look at existing
research and how some of these problems are dealt with and how some advantages are successfully exploited.

\paragraph{Proposing our own implementation}
After listing implementative and theoretical challenges we will proceed to propose our own implementation
of a federated learning system on a blockchain network.

We will try to compare it with the research papers discussed over the course of the document and provide our
own unique opinions and implementation choices ragarding topics with:
\begin{itemize}
  \item Outlining the network, consensus, and Blockchain layers to integrate the two
    technologies effectively.
  \item Exploring implementation choices, such as public versus private Blockchains,
    consensus mechanisms, and the role of smart contracts.
  \item Dealing with malicious updates in a robust manner.
  \item Making the system communication efficient and address storage and availability challenges.
\end{itemize}

We will then proceed to evaluate our implementation and compare it to a classic client-server implementation
and show the feasibility of blockchain integration.

\paragraph{Contributing to Academic research}
Saying that one of the aims of this document is to contribute to academic research may sound a bit pretentious
as this is by no means a groundbreaking work. However, by gathering existing research and providing our own
implementation and opinions on some relevant topics we hope that this document will be useful to someone
looking for a starting point in this field.

\paragraph{Conclusion}
By achieving these objectives, the thesis endeavors to shed light on the transformative
potential of combining Blockchain and Federated Learning, laying the groundwork for innovative
solutions to contemporary challenges in privacy-preserving and secure collaborative learning.


\subsection{Structure of the Document}

The remainder of this document is organized as follows:
\begin{itemize}
  \item In this chaper we will delve into the fundamentals of blockchain technology, examining its
    mechanisms, benefits, and limitations, focusing on the Bitcoin network as a case study.

  \item In chapter \ref{chap:federated-learning} we introduce federated learning. As a preliminary, we
    propose a very basic introduction to machine learning in \ref{sec:supervised-learning}, which
    introduces some topics and terminology relevant to federated learning. Then we proceed to explore the
    workings of the original federated learning proposal in \ref{sec:vanilla-fl}.
  \item In chapter \ref{chap:integration} we explore the possibility of integrating the two technologies.
    First we explore advantages, challenges and existing solutions in \ref{sec:blockchain-fl}. Then we
    proceed to present our implementation and our unique views on how to solve some challenges in
    \ref{chap:implementation}.
    We then evaluate our implementation against vanilla federated learning in \ref{chap:evaluation}.
  \item Finally we sum up the key findings and insights of this document in \ref{chap:conclusion}.
\end{itemize}


\clearpage{\pagestyle{empty}}
\section{Preliminaries}\label{chap:introduction}

\subsection{Context of the Thesis}\label{chap:contextandobjectives}
Advancements in technology have driven a data-centric shift across industries such as healthcare, finance,
smart cities, and the Internet of Things (IoT). This growing reliance on data has introduced significant
challenges in privacy, security, and scalability, necessitating innovative solutions. This thesis explores
the intersection of two emerging paradigms—Federated Learning (FL) and Blockchain technology—as a potential
means to address these challenges.

\textbf{Federated Learning} has emerged as a transformative approach to machine learning by enabling
collaborative model training without the need to share raw data. This is achieved by keeping relevant
computation local to devices and only sharing updates with the interested party.
This preserves data privacy while still achieving high-performance learning. Despite its benefits, FL faces
critical challenges, including reliance on a central server, susceptibility to malicious participants, and a
lack of transparency in the training process.

\textbf{Blockchain technology}, initially developed as the backbone for cryptocurrencies like Bitcoin,
provides a decentralized, secure, and transparent framework for managing data and operations among untrusted
parties. With its core features—such as consensus mechanisms, immutability, and smart contracts—Blockchain
has demonstrated value in sectors ranging from finance to supply chain management. However, its scalability
limitations and resource-intensive operations can hinder its applicability in dynamic environments.

Integrating FL and Blockchain presents a promising avenue for overcoming the limitations of both
technologies. Blockchain's decentralization can eliminate the dependency on a central server in FL, while its
transparency and consensus protocols enhance trust and accountability. Conversely, FL's distributed learning
capabilities reduce the need for raw data storage on the Blockchain, mitigating its scalability concerns.

This thesis is motivated by real-world challenges such as:
\begin{itemize}
  \item Ensuring privacy and security in data-sharing across regulated sectors like healthcare and finance.
  \item Establishing trust in collaborative environments where participants may not fully trust each other,
    such as in decentralized IoT systems.
\end{itemize}

While Federated Learning and Blockchain are an ongoing field of research and a fair amount of articles have
been published about these topics,
their integration has been underexplored. Combining the two technologies introduces novel challenges
including communication privacy and efficiency, synchronism of updates, security and scalability.

Through the design and evaluation of a Blockchain-enhanced Federated Learning framework, this thesis aims to
provide a foundational understanding of the two technologies and explore the challanges and opportunities in
context of privacy, trust, and scalability. The following section outlines the specific objectives of this
thesis, which are derived from this contextual foundation.

\subsection{Objectives of the Thesis}

This document aims to explore the feasibility and practicality of implementing a federated learning system on
a blockchain network. The contributions of this work are threefold:

\paragraph{Explore the technologies}
The thesis seeks to provide to the reader a foundational understanding of both Federated Learning and Blockchain
technologies. We will try to be as thorough as possible, within the limits of the purposes of this document,
highlighting the individual strengths, weaknesses, and potential areas of
application of both technologies.

\paragraph{Exploring the integration}
After providing an understanding of the technologies and their strengths and weaknesses we want to discuss
the advantages and challenges that arise when joining them.

We will discuss them in depth in the coming chapters but as an anticipation, some of the advantages include:
\begin{itemize}
  \item Addressing the central server reliance in FL through Blockchain's decentralized
    infrastructure.
  \item Enhancing trust, transparency, and security in FL systems using Blockchain's
    consensus mechanisms and immutability.
  \item Provide incentives for devices to partake in the federated learning process.
\end{itemize}

And some of the challenges include:
\begin{itemize}
  \item Computational and communication overhead due to Blockchain protocols.
  \item Ensuring security and consistence of the gloabl model.
\end{itemize}

When discussing the challenges and advantages we will also take a look at existing
research and how some of these problems are dealt with and how some advantages are successfully exploited.

\paragraph{Proposing our own implementation}
After listing implementative and theoretical challenges we will proceed to propose our own implementation
of a federated learning system on a blockchain network.

We will try to compare it with the research papers discussed over the course of the document and provide our
own unique opinions and implementation choices ragarding topics with:
\begin{itemize}
  \item Outlining the network, consensus, and Blockchain layers to integrate the two
    technologies effectively.
  \item Exploring implementation choices, such as public versus private Blockchains,
    consensus mechanisms, and the role of smart contracts.
  \item Dealing with malicious updates in a robust manner.
  \item Making the system communication efficient and address storage and availability challenges.
\end{itemize}

We will then proceed to evaluate our implementation and compare it to a classic client-server implementation
and show the feasibility of blockchain integration.

\paragraph{Contributing to Academic research}
Saying that one of the aims of this document is to contribute to academic research may sound a bit pretentious
as this is by no means a groundbreaking work. However, by gathering existing research and providing our own
implementation and opinions on some relevant topics we hope that this document will be useful to someone
looking for a starting point in this field.

\paragraph{Conclusion}
By achieving these objectives, the thesis endeavors to shed light on the transformative
potential of combining Blockchain and Federated Learning, laying the groundwork for innovative
solutions to contemporary challenges in privacy-preserving and secure collaborative learning.


\subsection{Structure of the Document}

The remainder of this document is organized as follows:
\begin{itemize}
  \item In this chaper we will delve into the fundamentals of blockchain technology, examining its
    mechanisms, benefits, and limitations, focusing on the Bitcoin network as a case study.

  \item In chapter \ref{chap:federated-learning} we introduce federated learning. As a preliminary, we
    propose a very basic introduction to machine learning in \ref{sec:supervised-learning}, which
    introduces some topics and terminology relevant to federated learning. Then we proceed to explore the
    workings of the original federated learning proposal in \ref{sec:vanilla-fl}.
  \item In chapter \ref{chap:integration} we explore the possibility of integrating the two technologies.
    First we explore advantages, challenges and existing solutions in \ref{sec:blockchain-fl}. Then we
    proceed to present our implementation and our unique views on how to solve some challenges in
    \ref{chap:implementation}.
    We then evaluate our implementation against vanilla federated learning in \ref{chap:evaluation}.
  \item Finally we sum up the key findings and insights of this document in \ref{chap:conclusion}.
\end{itemize}


\clearpage{\pagestyle{empty}}
\section{Preliminaries}\label{chap:introduction}

\subsection{Context of the Thesis}\label{chap:contextandobjectives}
Advancements in technology have driven a data-centric shift across industries such as healthcare, finance,
smart cities, and the Internet of Things (IoT). This growing reliance on data has introduced significant
challenges in privacy, security, and scalability, necessitating innovative solutions. This thesis explores
the intersection of two emerging paradigms—Federated Learning (FL) and Blockchain technology—as a potential
means to address these challenges.

\textbf{Federated Learning} has emerged as a transformative approach to machine learning by enabling
collaborative model training without the need to share raw data. This is achieved by keeping relevant
computation local to devices and only sharing updates with the interested party.
This preserves data privacy while still achieving high-performance learning. Despite its benefits, FL faces
critical challenges, including reliance on a central server, susceptibility to malicious participants, and a
lack of transparency in the training process.

\textbf{Blockchain technology}, initially developed as the backbone for cryptocurrencies like Bitcoin,
provides a decentralized, secure, and transparent framework for managing data and operations among untrusted
parties. With its core features—such as consensus mechanisms, immutability, and smart contracts—Blockchain
has demonstrated value in sectors ranging from finance to supply chain management. However, its scalability
limitations and resource-intensive operations can hinder its applicability in dynamic environments.

Integrating FL and Blockchain presents a promising avenue for overcoming the limitations of both
technologies. Blockchain's decentralization can eliminate the dependency on a central server in FL, while its
transparency and consensus protocols enhance trust and accountability. Conversely, FL's distributed learning
capabilities reduce the need for raw data storage on the Blockchain, mitigating its scalability concerns.

This thesis is motivated by real-world challenges such as:
\begin{itemize}
  \item Ensuring privacy and security in data-sharing across regulated sectors like healthcare and finance.
  \item Establishing trust in collaborative environments where participants may not fully trust each other,
    such as in decentralized IoT systems.
\end{itemize}

While Federated Learning and Blockchain are an ongoing field of research and a fair amount of articles have
been published about these topics,
their integration has been underexplored. Combining the two technologies introduces novel challenges
including communication privacy and efficiency, synchronism of updates, security and scalability.

Through the design and evaluation of a Blockchain-enhanced Federated Learning framework, this thesis aims to
provide a foundational understanding of the two technologies and explore the challanges and opportunities in
context of privacy, trust, and scalability. The following section outlines the specific objectives of this
thesis, which are derived from this contextual foundation.

\subsection{Objectives of the Thesis}

This document aims to explore the feasibility and practicality of implementing a federated learning system on
a blockchain network. The contributions of this work are threefold:

\paragraph{Explore the technologies}
The thesis seeks to provide to the reader a foundational understanding of both Federated Learning and Blockchain
technologies. We will try to be as thorough as possible, within the limits of the purposes of this document,
highlighting the individual strengths, weaknesses, and potential areas of
application of both technologies.

\paragraph{Exploring the integration}
After providing an understanding of the technologies and their strengths and weaknesses we want to discuss
the advantages and challenges that arise when joining them.

We will discuss them in depth in the coming chapters but as an anticipation, some of the advantages include:
\begin{itemize}
  \item Addressing the central server reliance in FL through Blockchain's decentralized
    infrastructure.
  \item Enhancing trust, transparency, and security in FL systems using Blockchain's
    consensus mechanisms and immutability.
  \item Provide incentives for devices to partake in the federated learning process.
\end{itemize}

And some of the challenges include:
\begin{itemize}
  \item Computational and communication overhead due to Blockchain protocols.
  \item Ensuring security and consistence of the gloabl model.
\end{itemize}

When discussing the challenges and advantages we will also take a look at existing
research and how some of these problems are dealt with and how some advantages are successfully exploited.

\paragraph{Proposing our own implementation}
After listing implementative and theoretical challenges we will proceed to propose our own implementation
of a federated learning system on a blockchain network.

We will try to compare it with the research papers discussed over the course of the document and provide our
own unique opinions and implementation choices ragarding topics with:
\begin{itemize}
  \item Outlining the network, consensus, and Blockchain layers to integrate the two
    technologies effectively.
  \item Exploring implementation choices, such as public versus private Blockchains,
    consensus mechanisms, and the role of smart contracts.
  \item Dealing with malicious updates in a robust manner.
  \item Making the system communication efficient and address storage and availability challenges.
\end{itemize}

We will then proceed to evaluate our implementation and compare it to a classic client-server implementation
and show the feasibility of blockchain integration.

\paragraph{Contributing to Academic research}
Saying that one of the aims of this document is to contribute to academic research may sound a bit pretentious
as this is by no means a groundbreaking work. However, by gathering existing research and providing our own
implementation and opinions on some relevant topics we hope that this document will be useful to someone
looking for a starting point in this field.

\paragraph{Conclusion}
By achieving these objectives, the thesis endeavors to shed light on the transformative
potential of combining Blockchain and Federated Learning, laying the groundwork for innovative
solutions to contemporary challenges in privacy-preserving and secure collaborative learning.


\subsection{Structure of the Document}

The remainder of this document is organized as follows:
\begin{itemize}
  \item In this chaper we will delve into the fundamentals of blockchain technology, examining its
    mechanisms, benefits, and limitations, focusing on the Bitcoin network as a case study.

  \item In chapter \ref{chap:federated-learning} we introduce federated learning. As a preliminary, we
    propose a very basic introduction to machine learning in \ref{sec:supervised-learning}, which
    introduces some topics and terminology relevant to federated learning. Then we proceed to explore the
    workings of the original federated learning proposal in \ref{sec:vanilla-fl}.
  \item In chapter \ref{chap:integration} we explore the possibility of integrating the two technologies.
    First we explore advantages, challenges and existing solutions in \ref{sec:blockchain-fl}. Then we
    proceed to present our implementation and our unique views on how to solve some challenges in
    \ref{chap:implementation}.
    We then evaluate our implementation against vanilla federated learning in \ref{chap:evaluation}.
  \item Finally we sum up the key findings and insights of this document in \ref{chap:conclusion}.
\end{itemize}


\clearpage{\pagestyle{empty}}
\section{Preliminaries}\label{chap:introduction}

\subsection{Context of the Thesis}\label{chap:contextandobjectives}
Advancements in technology have driven a data-centric shift across industries such as healthcare, finance,
smart cities, and the Internet of Things (IoT). This growing reliance on data has introduced significant
challenges in privacy, security, and scalability, necessitating innovative solutions. This thesis explores
the intersection of two emerging paradigms—Federated Learning (FL) and Blockchain technology—as a potential
means to address these challenges.

\textbf{Federated Learning} has emerged as a transformative approach to machine learning by enabling
collaborative model training without the need to share raw data. This is achieved by keeping relevant
computation local to devices and only sharing updates with the interested party.
This preserves data privacy while still achieving high-performance learning. Despite its benefits, FL faces
critical challenges, including reliance on a central server, susceptibility to malicious participants, and a
lack of transparency in the training process.

\textbf{Blockchain technology}, initially developed as the backbone for cryptocurrencies like Bitcoin,
provides a decentralized, secure, and transparent framework for managing data and operations among untrusted
parties. With its core features—such as consensus mechanisms, immutability, and smart contracts—Blockchain
has demonstrated value in sectors ranging from finance to supply chain management. However, its scalability
limitations and resource-intensive operations can hinder its applicability in dynamic environments.

Integrating FL and Blockchain presents a promising avenue for overcoming the limitations of both
technologies. Blockchain's decentralization can eliminate the dependency on a central server in FL, while its
transparency and consensus protocols enhance trust and accountability. Conversely, FL's distributed learning
capabilities reduce the need for raw data storage on the Blockchain, mitigating its scalability concerns.

This thesis is motivated by real-world challenges such as:
\begin{itemize}
  \item Ensuring privacy and security in data-sharing across regulated sectors like healthcare and finance.
  \item Establishing trust in collaborative environments where participants may not fully trust each other,
    such as in decentralized IoT systems.
\end{itemize}

While Federated Learning and Blockchain are an ongoing field of research and a fair amount of articles have
been published about these topics,
their integration has been underexplored. Combining the two technologies introduces novel challenges
including communication privacy and efficiency, synchronism of updates, security and scalability.

Through the design and evaluation of a Blockchain-enhanced Federated Learning framework, this thesis aims to
provide a foundational understanding of the two technologies and explore the challanges and opportunities in
context of privacy, trust, and scalability. The following section outlines the specific objectives of this
thesis, which are derived from this contextual foundation.

\subsection{Objectives of the Thesis}

This document aims to explore the feasibility and practicality of implementing a federated learning system on
a blockchain network. The contributions of this work are threefold:

\paragraph{Explore the technologies}
The thesis seeks to provide to the reader a foundational understanding of both Federated Learning and Blockchain
technologies. We will try to be as thorough as possible, within the limits of the purposes of this document,
highlighting the individual strengths, weaknesses, and potential areas of
application of both technologies.

\paragraph{Exploring the integration}
After providing an understanding of the technologies and their strengths and weaknesses we want to discuss
the advantages and challenges that arise when joining them.

We will discuss them in depth in the coming chapters but as an anticipation, some of the advantages include:
\begin{itemize}
  \item Addressing the central server reliance in FL through Blockchain's decentralized
    infrastructure.
  \item Enhancing trust, transparency, and security in FL systems using Blockchain's
    consensus mechanisms and immutability.
  \item Provide incentives for devices to partake in the federated learning process.
\end{itemize}

And some of the challenges include:
\begin{itemize}
  \item Computational and communication overhead due to Blockchain protocols.
  \item Ensuring security and consistence of the gloabl model.
\end{itemize}

When discussing the challenges and advantages we will also take a look at existing
research and how some of these problems are dealt with and how some advantages are successfully exploited.

\paragraph{Proposing our own implementation}
After listing implementative and theoretical challenges we will proceed to propose our own implementation
of a federated learning system on a blockchain network.

We will try to compare it with the research papers discussed over the course of the document and provide our
own unique opinions and implementation choices ragarding topics with:
\begin{itemize}
  \item Outlining the network, consensus, and Blockchain layers to integrate the two
    technologies effectively.
  \item Exploring implementation choices, such as public versus private Blockchains,
    consensus mechanisms, and the role of smart contracts.
  \item Dealing with malicious updates in a robust manner.
  \item Making the system communication efficient and address storage and availability challenges.
\end{itemize}

We will then proceed to evaluate our implementation and compare it to a classic client-server implementation
and show the feasibility of blockchain integration.

\paragraph{Contributing to Academic research}
Saying that one of the aims of this document is to contribute to academic research may sound a bit pretentious
as this is by no means a groundbreaking work. However, by gathering existing research and providing our own
implementation and opinions on some relevant topics we hope that this document will be useful to someone
looking for a starting point in this field.

\paragraph{Conclusion}
By achieving these objectives, the thesis endeavors to shed light on the transformative
potential of combining Blockchain and Federated Learning, laying the groundwork for innovative
solutions to contemporary challenges in privacy-preserving and secure collaborative learning.


\subsection{Structure of the Document}

The remainder of this document is organized as follows:
\begin{itemize}
  \item In this chaper we will delve into the fundamentals of blockchain technology, examining its
    mechanisms, benefits, and limitations, focusing on the Bitcoin network as a case study.

  \item In chapter \ref{chap:federated-learning} we introduce federated learning. As a preliminary, we
    propose a very basic introduction to machine learning in \ref{sec:supervised-learning}, which
    introduces some topics and terminology relevant to federated learning. Then we proceed to explore the
    workings of the original federated learning proposal in \ref{sec:vanilla-fl}.
  \item In chapter \ref{chap:integration} we explore the possibility of integrating the two technologies.
    First we explore advantages, challenges and existing solutions in \ref{sec:blockchain-fl}. Then we
    proceed to present our implementation and our unique views on how to solve some challenges in
    \ref{chap:implementation}.
    We then evaluate our implementation against vanilla federated learning in \ref{chap:evaluation}.
  \item Finally we sum up the key findings and insights of this document in \ref{chap:conclusion}.
\end{itemize}


% BIBLIOGRAFIA
\nocite{*}
\addcontentsline{toc}{chapter}{Bibliography}
\label{Bibliography}
\bibliographystyle{IEEEtran}
\bibliography{bibliography}
\rhead[\fancyplain{}{\bfseries Bibliography}]
{\fancyplain{}{\bfseries\thepage}}
\end{document}
